\documentclass[11pt,a4paper]{amsart}

\usepackage[margin=1in]{geometry}
\usepackage[T1]{fontenc}
\usepackage[utf8]{inputenc}
\usepackage{lmodern}
\usepackage{microtype}
\usepackage{amsmath,amssymb,amsthm,mathtools}
\usepackage{enumitem}
\usepackage[colorlinks=true,linkcolor=blue,citecolor=blue,urlcolor=blue]{hyperref}
\usepackage[capitalize,noabbrev]{cleveref}

\numberwithin{equation}{section}

\newtheorem{theorem}{Theorem}[section]
\newtheorem{proposition}[theorem]{Proposition}
\newtheorem{lemma}[theorem]{Lemma}
\newtheorem{corollary}[theorem]{Corollary}
\newtheorem{definition}[theorem]{Definition}
\theoremstyle{remark}
\newtheorem{remark}[theorem]{Remark}

\DeclareMathOperator{\RS}{RS}
\DeclareMathOperator{\Span}{Span}
\DeclareMathOperator{\Conf}{Conf}
\DeclareMathOperator{\supp}{supp}

\newcommand{\F}{\mathbb F}
\newcommand{\B}{\mathbb B}
\newcommand{\Z}{\mathbb Z}
\newcommand{\Ccal}{\mathcal C}
\newcommand{\Lcal}{\mathcal L}
\newcommand{\Scal}{\mathcal S}
\newcommand{\Wcal}{\mathcal W}
\newcommand{\OmegaC}{\Omega^{\circ}}
\newcommand{\Gord}{\Gamma^{\rm ord}}
\newcommand{\Gsid}{\Gamma^{\rm sid}}
\newcommand{\barN}{\overline N}
\newcommand{\Eprof}{\mathfrak E}
\newcommand{\eps}{\varepsilon}
\newcommand{\abs}[1]{\lvert #1\rvert}
\newcommand{\defeq}{\coloneqq}

\title{Asymptotic RS MCA}
\date{}
\hypersetup{pdftitle={Asymptotic RS MCA}}

\begin{document}

\begin{abstract}
We give a conditional asymptotic Reed--Solomon mutual-correlated-agreement compiler from a first-match profile ledger.  After quotient, planted, and other algebraic profiles have been routed, Balog--Szemeredi--Gowers and quasicube growth eliminate the high-energy branch of the primitive prefix problem, provided its Sidon-heavy branch is paid separately.  We construct smooth rows with \(\log|\B_n|=O(\log n)=o(n)\) having an exponentially large Sidon-heavy square-quotient prefix fiber at the identity crossing and, after scalar extension, exponentially many distinct MCA-bad slopes.  The construction persists after stabilizer deletion, while its periodic form transfers to standard circle codes.  Consequently the frontier theorem must use the full profile envelope, a Sidon payment on the routed primitive residual, and a distinct-ray compiler; the identity formula is recovered only under explicit identity-dominance and target-reserve hypotheses.
\end{abstract}

\maketitle

\section{Statement}

Let \(D\subseteq \F\) be an evaluation set of size \(n\), and let
\[
        C=\RS_{\F}(D,k)=\{(f(t))_{t\in D}:\deg f<k\}.
\]
For \(a\in\{0,\ldots,n\}\), a finite slope \(\gamma\in\F\) is support-wise MCA-bad for a pair \(r_1,r_2\in\F^D\) at agreement \(a\) if some set \(S\subseteq D\), \(|S|\ge a\), and some \(h\in\F[X]\), \(\deg h<k\), satisfy \(r_1+\gamma r_2=h\) on \(S\), while \(r_1,r_2\) are not simultaneously explained on the same \(S\) by two degree-\(<k\) polynomials.  Put
\[
        B_C^{\rm MCA}(a)
        =\max_{r_1,r_2}\#\{\gamma\in\F:\gamma\text{ is MCA-bad at agreement }a\}.
\]
We use the full slope field, so \(q_{\rm line}=|\F|\), and set
\(\eps_{\rm mca}(C,1-a/n)=B_C^{\rm MCA}(a)/|\F|\).  Related finite-row
normalizations appear in \cite{Cho26CapV13,Cho26Grande}.

We consider smooth multiplicative or circle-type RS row sequences
\(C_n=\RS_{\F_n}(D_n,k_n)\) with \(|D_n|=n\),
\(\rho_n=k_n/n\to\rho\in(0,1)\), and
\(D_n\subseteq\B_n\subseteq\F_n\), where
\(\beta_n=\log_2|\B_n|=o(n)\).  For an agreement
\(a_n=k_n+1+w_n\), define the identity-prefix average
\[
        \barN_{n,a_n}=\binom n{a_n}|\B_n|^{-w_n}.
\]
For \(r\in(0,1)\) and \(b\ge0\), put
\[
        g^*(r,b)=\sup\{g\in[0,1-r]:H_2(r+g)\ge bg\},
\]
where \(H_2(x)=-x\log_2x-(1-x)\log_2(1-x)\), and write
\(g_n^*=g^*(\rho_n,\beta_n)\).

Here a multiplicatively smooth domain is a coset \(\theta H\) of a
cyclic subgroup of \(\B_n^\times\), together with its uniform-fiber power
quotients.  A circle-type domain is a Chebyshev image of a twin coset
\(gH\cup g^{-1}H\) in a norm-one torus, or a circle evaluation set
obtained from such a multiplicative set by diagonal uniformization or
stereographic parametrization.  More generally, a polynomial map
\(\varphi\) of degree \(d\) is a smooth quotient of \(D\) when every
nonempty fiber \(\varphi^{-1}(y)\cap D\) has exactly \(d\) elements.  The
quantitative assumptions used below are the stated profile-ledger and ray
conditions; these geometric names specify the profile atlas to which they
are applied.

For a fixed row and agreement, a \emph{first-match profile atlas} is fixed
before the received line is chosen.  It is an ordered finite collection of
constructible witness charts.  A profile \(\lambda\) in the atlas specifies
a full support slice
\(\Omega^0_\lambda\), a coefficient field \(\B_\lambda\), and
\(R_\lambda\) effective prefix equations; its natural average is
\(\barN_\lambda=|\Omega^0_\lambda||\B_\lambda|^{-R_\lambda}\).  For a
received line, \(\Lambda(r_1,r_2;a)\) is the subcollection of fixed atlas
profiles that occur under first match.  Exhaustiveness of this atlas is an
explicit ledger hypothesis, so profiles cannot be added or deleted after
their sizes are known.  The identity profile is always included.  Define
the profile envelope
\[
 \Eprof_n(a)=1+(n-a+1)+
 \sup_{r_1,r_2}\sum_{\lambda\in\Lambda(r_1,r_2;a)}\barN_\lambda.
 \tag{1.1}\label{eq:profile-envelope}
\]
The tangent term is displayed because it cannot be absorbed when the
identity average is below one, and
\(\Eprof_n(a)\ge1+(n-a+1)+\barN_{n,a}\).  A window is
\emph{identity-dominant} if some \(\kappa_n=\exp(o(n))\), uniformly in the
window, satisfies
\[
 \Eprof_n(a)\le
 \kappa_n\bigl(1+(n-a+1)+\barN_{n,a}\bigr).
\]
On the agreement lattice the target threshold is
\[
 \delta_C^*(\eps)=\sup\left\{1-\frac an:
 a\in\{k+1,\ldots,n\},\ \eps_{\rm mca}(C,1-a/n)\le\eps\right\},
\]
whenever this set is nonempty.

\begin{definition}[Cells and first-match payment]\label{def:cells}
A \emph{cell} is a constructible subfamily of the MCA witness incidence
space, or a bounded-complexity family of such subfamilies, stable under
extension of the ground field.  Cells are ordered.  A slope is assigned to
the first cell containing one of its witnesses.  A nonprimitive algebraic
cell is \emph{paid} at agreement \(a_n\) if its distinct-slope image is at
most its printed natural-profile budget; a Fourier/Sidon cell is paid in
the support statistic of \cref{def:sidon-paid}.  A profile family is paid
if the sum of its first-match budgets is at most
\(\exp(o(n))\Eprof_n(a_n)\).  Only in an identity-dominant window may this
be specialized to the identity budget.
\end{definition}

\begin{lemma}[First-match disjointization]\label{lem:first-match}
If ordered cells \(\Ccal_1,\ldots,\Ccal_J\) cover a set of bad slopes and the \(j\)-th first-match cell has budget \(U_j\), then the covered slopes contribute at most \(\sum_jU_j\) to \(B_C^{\rm MCA}(a)\).
\end{lemma}

\begin{proof}
For a fixed received line, assign each bad slope to the least-indexed cell containing a witness.  These assigned slope classes are disjoint, and the \(j\)-th class is contained in the projection of \(\Ccal_j\) after earlier cells have been removed.  Summing the budgets and then maximizing over received lines proves the claim.
\end{proof}

\begin{definition}[Closed paid ledger]\label{def:closed-ledger}
A smooth or circle-type row sequence is \emph{closed-ledger} at agreements
\(a_n\) if
the following first-match cells cover all MCA-bad witnesses except the
primitive leaves of \cref{def:primitive-leaf}, the number of profiles is
\(\exp(o(n))\), their natural coefficient fields and prefix depths are
retained in \eqref{eq:profile-envelope}, the algebraic cells have total
distinct-slope contribution \(\exp(o(n))\Eprof_n(a_n)\), and the
Fourier/Sidon and primitive support statistics are paid at their declared
profile scales.
\begin{enumerate}[label=\textup{(C\arabic*)},leftmargin=2.7em]
    \item \emph{Quotient-pullback cells.}  The support or locator descends along a nontrivial finite map \(D\to D'\).  These cells include quotient-remainder locators, divisor-block support ledgers, quotient-image lcm ledgers, and quotient pullback recursion.
    \item \emph{Chebyshev and dihedral cells.}  On multiplicative or circle rows, the support is invariant under the dihedral action, or the locator factors through a Chebyshev or circle quotient such as \(z+z^{-1}\).
    \item \emph{Planted-block cells.}  A positive-density block is forced into the support or common zero set.  The payment is the entropy cost of planting that block, together with the bounded quotient profile.
    \item \emph{Tangent and deep-center cells.}  The received line lies in a tangent, secant, or high-contact stratum of the RS code variety, including the high-agreement tangent cell.
    \item \emph{Extension and descent cells.}  The witness either descends to a smaller field or is controlled by a bounded extension-line or extension-pole parameter.
    \item \emph{Differential-locator and regular-rank cells.}  The locator map has a rank drop, a ramification defect, or an overdetermined Hankel minor.  These cells include regular closed-range Hankel packing, canonical rank-drop ledgers, and curve rank-drop ledgers.
    \item \emph{Saturation and effective-image-collapse cells.}  The apparent support count is larger than the actual line-ray or codeword-ray image.  The payment uses exact saturation identities and exact residual image ledgers.
    \item \emph{Balanced-core and split-pencil cells.}  The residual witness lies in a balanced-core split-pencil chart.  One-parameter primitive locator pencils are paid by the moving-root theorem; deficiency-one and identically split top charts are tangent-borne or eliminated.  Every higher-dimensional residual chart requires a direct distinct-ray estimate.
    \item \emph{Fourier/Sidon cells.}  In a primitive prefix leaf, a heavy fiber with exponentially small additive closure is cut as a Fourier/Sidon cell: major arcs are routed to \textup{(C1)--(C8)}, and minor arcs are paid by Fourier flatness.
\end{enumerate}
The remaining first-match complement is a disjoint union of primitive
prefix leaves.  Every such leaf has the Boolean fixed-density Vandermonde
form of \cref{def:primitive-leaf}.  If it is attached to profile
\(\lambda\), with residual mass \(M_\lambda\), residual image size
\(L_\lambda\), and active length \(N_\lambda\), closure requires
\[
 \log(M_\lambda/L_\lambda)=o(N_\lambda),\qquad
 \frac{M_\lambda}{L_\lambda}
 \le\exp(o(n))\barN_\lambda.
\]
The \(o\)-terms are uniform over the atlas.  Its Sidon payment is imposed
only after all earlier profiles have been
removed.  Thus a closed profile ledger includes the residual-to-full
average comparison.  Conversion of the resulting support budgets into
distinct MCA slopes is the separate RC condition of
\cref{def:ray-compiler}.
\end{definition}

\begin{definition}[Direct ray compiler]\label{def:ray-compiler}
A row satisfies \emph{RC} at agreement \(a_n\) if the distinct slopes
projected from every primitive residual profile are at most
\(\exp(o(n))\) times that profile's paid budget, with total cost
\(\exp(o(n))\Eprof_n(a_n)\).  A max-fiber or support-pair estimate alone
is not RC: it counts supports or pairs, not the image of the witness
incidence in slope space.
\end{definition}

\begin{theorem}[Cellwise ledger interface]\label{thm:closed-ledger-package}
The results cited in \cite{Cho26CapV13,Cho26Grande} supply the cellwise
constructions and estimates in \textup{(C1)--(C8)} and the Fourier-flat
part of \textup{(C9)}.  For a given smooth or circle row they assemble to
a closed profile ledger only when the atlas is exhaustive, every profile is
charged at its natural scale in \eqref{eq:profile-envelope}, and the
residual-to-full average comparison holds, and the Sidon-heavy part of the
routed primitive residual is paid.  The full MCA compiler additionally
requires RC.  None of these cellwise results implies identity dominance.
\end{theorem}

\begin{proof}
The cited cellwise estimates are as follows.

For \textup{(C1)}, the quotient-remainder lower and support ledgers, divisor-union support ledger, exact finite-parameter quotient-image ledger, affine/projective lcm ledgers, quotient pullback recursion, and bounded quotient census are proved in \cite[quotient-remainder, quotient support, quotient-image, and bounded-census ledgers]{Cho26CapV13}.  The corresponding first-match convention and quotient status cell are isolated in \cite{Cho26Grande}.

For \textup{(C2)}, exact Chebyshev fibers on circle twin-coset domains,
circle quotient transport, and circle profile estimates are proved in
\cite[Chebyshev and circle sections]{Cho26CapV13}.  Dihedral and Chebyshev
cells are routed as quotient cells with extra monodromy; their payment still
uses the natural-profile comparison required in the theorem.

For \textup{(C3)}, planted quotient profiles, planted quotient-core lower
counts, planted list-budget windows, and bounded active quotient order are
established in \cite[planted quotient profile and bounded quotient
census]{Cho26CapV13}.  They identify the relevant planted entropy scale;
the corresponding upper payment remains part of the closed-ledger
hypothesis.

For \textup{(C4)}, the exact high-agreement tangent cell and tangent-borne split-kernel sections are proved in \cite[high-agreement tangent exact cell and split-kernel tangent-borne theorem]{Cho26CapV13} and refined in \cite[exact high-agreement tangent cell]{Cho26Grande}.

For \textup{(C5)}, the extension-pole deep-list and quotient-remainder floors, extension-line dimension-degree ledger, subfield confinement, and genuinely extension-valued line accounting are proved in \cite[extension-pole and extension-line ledgers]{Cho26CapV13}; the full-extension cell target and extension-valued rank-one floor are recorded in \cite{Cho26Grande}.

For \textup{(C6)}, regular Hankel minors, canonical exact-agreement rank-drop ledgers, closed-ball rank-drop lcm ledgers, curve rank-drop ledgers, arbitrary-residual image ledgers, and the scanner residual ledger are proved in \cite[Hankel, rank-drop, curve, residual-image, and scanner ledgers]{Cho26CapV13}.  These are exactly the differential-locator and regular-rank payments.

For \textup{(C7)}, the exact saturation identity and line-ray saturation identity are proved in \cite[saturation and line-ray saturation identities]{Cho26Grande}; exact arbitrary-residual image and closed-ball residual image ledgers appear in \cite[residual image ledgers]{Cho26CapV13}.  Thus raw support overcounts are replaced by actual image counts.

For \textup{(C8)}, the interpolation-lattice split-pencil reduction, near-rational payment, deficiency-one split-test eliminant, split annihilator dictionary, tangent-borne split charts, boundary-profile identification with Q, base-field split-pencil floor, one-parameter moving-root theorem, and primitive one-pencil corollary are proved in \cite[split-pencil and BC sections]{Cho26Grande}, with the balanced-core and split-pencil reductions also developed in \cite[rank-one and split-pencil reductions]{Cho26CapV13}.

For \textup{(C9)}, Fourier-flat leaves satisfy asymptotic Q and
large-characteristic Fourier examples are proved in
\cite[Fourier-flat Q section]{Cho26Grande}.  The Sidon cut and its
high-energy alternative are stated and proved below.  Fourier flatness pays
the minor-arc branch after algebraic major arcs have been routed; it does
not give a uniform identity-normalized payment before quotient, subfield,
planted, and remainder profiles are removed.

Finally, first-match disjointization is \cref{lem:first-match}.  These
results give the asserted cellwise interface.  Exhaustion, comparison at
the natural profile scales, residual-to-full average comparison, residual
Sidon payment, and projection from support statistics to distinct slopes
are precisely the additional conditions in the statement.
\end{proof}

Thus the cited geometry identifies and estimates the declared cells but
does not unconditionally close an identity-scale ledger.  The argument
below proves the analytic primitive implication and the target compiler
once the printed profile, Sidon, and ray conditions have been verified.

\begin{theorem}[Conditional profile-envelope frontier]\label{thm:frontier}
Let \(C_n=\RS_{\F_n}(D_n,k_n)\) be a smooth or circle-type row as above,
let \(0\le\eps_n\le1\), put \(T_n=\lfloor\eps_n|\F_n|\rfloor\), and let
\(\kappa_n=\exp(o(n))\) dominate the compiler losses uniformly in the
window under consideration.  Define the first safe agreement
\[
 a_n^*=\min\{a\in\{k_n+1,\ldots,n\}:
                 B_{C_n}^{\rm MCA}(a)\le T_n\},
 \qquad \delta_n^*=1-a_n^*/n,
\]
whenever the set is nonempty.  Suppose agreements
\(a_{-,n}<a_{+,n}\) satisfy:
\begin{enumerate}[label=\textup{(\roman*)},leftmargin=2.4em]
\item a closed profile ledger and RC hold at \(a_{+,n}\), and
      \(\kappa_n\Eprof_n(a_{+,n})\le T_n\);
\item a certified identity, quotient, planted, or remainder construction
      gives \(U_n(a_{-,n})\) distinct MCA-bad slopes, where
      \(U_n(a_{-,n})>T_n\).
\end{enumerate}
Then
\[
 a_{-,n}<a_n^*\le a_{+,n},\qquad
 1-\frac{a_{+,n}}n\le\delta_n^*<1-\frac{a_{-,n}}n.
 \tag{1.2}\label{eq:threshold-bracket}
\]
If the two agreements are \(o(n)\) apart around
\(a_n=k_n+1+g_n n\), then
\(\delta_n^*=1-\rho_n-g_n+o(1)\).  For the identity construction, when
\(|\F_n|>n\), one may take
\[
 U_n(a)=
 \left\lceil
 \frac{L_n(a)(|\F_n|-n)}
      {|\F_n|-n+k_n(L_n(a)-1)}
 \right\rceil,\qquad
 L_n(a)=
 \left\lceil\binom na|\B_n|^{-(a-k_n-1)}\right\rceil .
\]

For the identity-dominant specialization put
\[
 \tau_n=\frac1n\log_2(1+T_n),\qquad
 F_n(g)=H_2(\rho_n+g)-\beta_ng,
\]
and define the target-adjusted right crossing by
\[
 g_{T,n}=\sup\bigl(\{g\in[0,1-\rho_n]:F_n(g)\ge\tau_n\}\cup\{0\}\bigr).
 \tag{1.3}\label{eq:target-crossing}
\]
At an interior crossing,
\[
 H_2(\rho_n+g_{T,n})-\beta_ng_{T,n}
   =\frac1n\log_2(1+T_n).
\]
If identity dominance and the collision-aware lower construction supply
agreements satisfying \textup{(i)--(ii)} with
\(a_{\pm,n}=k_n+1+g_{T,n}n+o(n)\), then
\(\delta_n^*=1-\rho_n-g_{T,n}+o(1)\).  If additionally
\(T_n=\exp(o(n))\), then \(g_{T,n}=g^*(\rho_n,\beta_n)+o(1)\), and hence
\[
 \delta_n^*=1-\rho_n-g^*(\rho_n,\beta_n)+o(1).
 \tag{1.4}\label{eq:identity-frontier}
\]
\end{theorem}

The rest of the paper proves this conditional implication from the
cellwise interface and the external additive-combinatorics facts stated in
\cref{sec:bsg}.  The counterexample below shows that identity dominance and
any pre-routing uniform Sidon payment cannot be assumed.  Residual Sidon
payment remains a separate hypothesis, while
\cref{rem:q-sp-no-ray} explains independently why RC is required.

\section{Primitive leaves}

\begin{definition}[Primitive prefix leaf]\label{def:primitive-leaf}
A primitive leaf consists of a finite set \(T\), \(|T|=N\), an integer \(m\) with \(m/N\) bounded away from \(0\) and \(1\), a field \(K\), nonzero weights \(\rho(t)\in K^\times\), and columns
\[
        v_t=\rho(t)(1,t,\ldots,t^{R-1})\in K^R,\qquad t\in T.
\]
The active support family is a first-match residual slice
\[
        \OmegaC\subseteq \{x\in\{0,1\}^{T}:\sum_{t\in T}x_t=m\},
\]
and the prefix map is \(\Phi(x)=\sum_{t\in T}x_tv_t\).  Write \(\Scal=\operatorname{im}\Phi\), \(L=|\Scal|\), \(M=|\OmegaC|\), \(\barN=M/L\), and \(F_s=\OmegaC\cap\Phi^{-1}(s)\).  The leaf is at the frontier if \(\log\barN=o(N)\).
\end{definition}

For \(q=q_N\to\infty\), define
\[
        \Gord_q=L^{-1}\sum_{s\in\Scal}\left(\frac{|F_s|}{\barN}\right)^q .
\]
We call \(q\) \emph{logarithmic} if
\(q_N\le(\log N)^A\) for some fixed \(A\); since every image fiber is
nonempty, \(L\le M\le2^N\), so \(\log L=o(Nq)\) for every such \(q\).

\begin{lemma}[Moment-max equivalence]\label{lem:moment-max}
If \(q\to\infty\) and \(\log L=o(Nq)\), then
\[
        L^{-1}\left(\max_s\frac{|F_s|}{\barN}\right)^q
        \le \Gord_q
        \le
        \left(\max_s\frac{|F_s|}{\barN}\right)^q .
\]
Hence \(\Gord_q\le \exp(o(Nq))\) is equivalent to \(\max_s|F_s|\le\exp(o(N))\barN\).
\end{lemma}

\begin{proof}
The upper bound replaces every summand by the maximum; the lower bound keeps a maximum summand.  Taking \(q\)-th roots and using \(\log L=o(Nq)\) gives the final assertion.
\end{proof}

\section{Sidon cut}

For a finite \(F\) in an abelian group, let
\[
        E(F)=\#\{(a,b,c,d)\in F^4:a-b=c-d\},\qquad
        \Delta(F)=E(F)/|F|^3.
\]
Thus \(\Delta(F)\) is the probability that \(a-b+c\in F\) for independent uniform \(a,b,c\in F\).
Whenever supports are used as an additive set, we identify
\(S\subseteq D\) with its incidence vector \(\mathbf1_S\in\Z^D\);
all Boolean energies therefore lie in this torsion-free ambient group.

For \(\sigma>0\), define the Sidon-heavy moment of a primitive leaf by
\[
        \Gsid_{q,\sigma}
        =L^{-1}\sum_{\Delta(F_s)\le e^{-\sigma N}}
        \left(\frac{|F_s|}{\barN}\right)^q .
\]

\begin{definition}[Paid Fourier/Sidon cell]\label{def:sidon-paid}
The Fourier/Sidon cell is paid if, for every fixed \(\sigma>0\) and every logarithmic \(q\), one has \(\Gsid_{q,\sigma}\le\exp(o(Nq))\).  Equivalently, no positive-rate heavy fiber with exponentially small additive closure remains in the primitive residual.
\end{definition}

In practice this cell is paid by Fourier flatness after major arcs have been routed to algebraic cells.  If \(\mu=\Phi_*\operatorname{Unif}(\OmegaC)\) and \(\sum_{\chi\ne1}|\widehat\mu(\chi)|\le\exp(o(N))\), Fourier inversion gives \(L\mu(s)\le\exp(o(N))\) for all \(s\), hence primitive Q.  Verification of a closed ledger must separately prove that every large Fourier coefficient is a quotient, descent, dihedral, ramification, or effective-image-collapse major arc; once verified, first match assigns it to the earlier cells in \cref{def:closed-ledger}.

\begin{proposition}[Energy extraction after the Sidon cut]\label{prop:energy-extract}
Assume the Fourier/Sidon cell is paid.  If \(\Gord_q\ge e^{\eta Nq}\) for some fixed \(\eta>0\), then some level \(s=s_N\) satisfies \(|F_s|\ge e^{\eta N-o(N)}\barN\) and \(E(F_s)\ge |F_s|^3e^{-o(N)}\).
\end{proposition}

\begin{proof}
For a fixed \(\sigma>0\), split \(\Gord_q\) into terms with
\(\Delta(F_s)\le e^{-\sigma N}\) and the complementary terms.  The first
part is \(\Gsid_{q,\sigma}=\exp(o(Nq))\), so a complementary term has
\((|F_s|/\barN)^q\ge e^{\eta Nq-o(Nq)}\).  Hence
\(|F_s|\ge e^{\eta N-o(N)}\barN\) and
\(\Delta(F_s)>e^{-\sigma N}\).

For completeness, apply this assertion successively with
\(\sigma_j=1/j\).  Along the subsequence on which the moment excess holds,
choose the \(j\)-th index after the threshold at which the paid estimate
for \(\sigma_j\) is valid.  The resulting diagonal subsequence has
\(\sigma_j\downarrow0\), and therefore
\(E(F_s)\ge |F_s|^3e^{-o(N)}\).
\end{proof}

\section{Boolean additive combinatorics}\label{sec:bsg}

We use the energy form of Balog--Szemeredi--Gowers \cite{BalogSzemeredi,Gowers1998}.  Its precise polynomial exponents are irrelevant here.

\begin{theorem}[Balog--Szemeredi--Gowers]\label{thm:bsg}
If \(A\) is a finite subset of an abelian group and \(E(A)\ge |A|^3/K\), then some \(A'\subseteq A\) satisfies \(|A'|\ge K^{-C}|A|\) and \(|A'-A'|\le K^C|A'|\), where \(C\) is absolute.
\end{theorem}

We also use the quasicube sumset theorem of Matolcsi--Ruzsa--Shakan--Zhelezov \cite{MRSZ2020} in the sharpened form of Green--Matolcsi--Ruzsa--Shakan--Zhelezov \cite{GMRSZ2020}: if \(U\subseteq\{0,1\}^d\), then \(|P+Q+U|\ge |P|^{1/2}|Q|^{1/2}|U|\) for finite nonempty \(P,Q\subseteq\Z^d\).

\begin{theorem}[Quasicube difference growth]\label{thm:quasicube}
Every finite \(A\subseteq\{0,1\}^{N}\) satisfies \(|A-A|\ge |A|^{3/2}\).
\end{theorem}

\begin{proof}
Apply the quasicube theorem with \(U=A\), \(P=-A\), and \(Q=\{0\}\).
\end{proof}

\begin{proposition}[No large high-energy Boolean fiber]\label{prop:no-high-energy}
There is no sequence \(F_N\subseteq\{0,1\}^{N}\) with \(|F_N|\ge e^{cN-o(N)}\) and \(E(F_N)\ge |F_N|^3e^{-o(N)}\) for fixed \(c>0\).
\end{proposition}

\begin{proof}
By \cref{thm:bsg} with \(K=e^{o(N)}\), a subset \(F'_N\subseteq F_N\) has \(|F'_N|\ge e^{cN-o(N)}\) and \(|F'_N-F'_N|\le e^{o(N)}|F'_N|\).  But \cref{thm:quasicube} gives \(|F'_N-F'_N|\ge |F'_N|^{3/2}\), so \(|F'_N|^{1/2}\le e^{o(N)}\), a contradiction.
\end{proof}

\begin{corollary}[Quantitative Boolean energy bound]
\label{cor:quantitative-boolean-energy}
There is an absolute constant \(A>0\) such that, for every
\(F\subseteq\{0,1\}^N\) and \(K\ge1\),
\[
 E(F)\ge |F|^3/K\qquad\Longrightarrow\qquad |F|\le K^A.
\]
\end{corollary}

\begin{proof}
Retain the polynomial dependence in \cref{thm:bsg}.  It gives
\(F'\subseteq F\) with \(|F'|\ge K^{-C}|F|\) and
\(|F'-F'|\le K^C|F'|\).  By \cref{thm:quasicube},
\(|F'|^{3/2}\le K^C|F'|\), hence \(|F'|\le K^{2C}\) and
\(|F|\le K^{3C}\).  Take \(A=3C\).
\end{proof}

\begin{theorem}[Primitive Q after the Sidon cut]\label{thm:primitive-q}
In every frontier primitive leaf with paid Fourier/Sidon cell, \(\max_s|F_s|\le\exp(o(N))\barN\).
\end{theorem}

\begin{proof}
If not, \cref{lem:moment-max} applies because
\(L\le M\le2^N\), and gives \(\Gord_q\ge e^{\eta Nq}\) for some fixed
\(\eta>0\) along a subsequence and some logarithmic \(q\).
\Cref{prop:energy-extract} then gives a fiber
\(F_s\subseteq\{0,1\}^{T}\) with
\(|F_s|\ge e^{\eta N-o(N)}\barN=e^{cN-o(N)}\) and
\(E(F_s)\ge |F_s|^3e^{-o(N)}\), because \(\log\barN=o(N)\).
This contradicts \cref{prop:no-high-energy}.
\end{proof}

\section{A polynomial-field obstruction}\label{sec:quotient-obstruction}

For \(S\in\binom Da\), write
\[
 Q_S(X)=\prod_{t\in S}(X-t)
       =X^a+q_1(S)X^{a-1}+\cdots+q_a(S),
 \qquad \Phi_w(S)=(q_1(S),\ldots,q_w(S)),
\]
and put \(\mathcal F_z=\Phi_w^{-1}(z)\) and
\(\barN_1=\binom na|\B|^{-w}\).  To distinguish the unrestricted identity
prefix from a routed primitive leaf, define
\[
 \widetilde\Gsid_{r,\sigma}
 =|\B|^{-w}\!\sum_{\substack{z\in\B^w:\ \mathcal F_z\ne\varnothing\\
                         \Delta(\mathcal F_z)\le e^{-\sigma n}}}
       \left(\frac{|\mathcal F_z|}{\barN_1}\right)^r .
 \tag{5.1}\label{eq:unrestricted-sidon-moment}
\]

\begin{theorem}[Smooth quotient, Sidon, and MCA obstruction]
\label{thm:polynomial-obstruction}
Fix \(0<\alpha<1/2\).  Let \(p\to\infty\) through odd primes, put
\[
 n=2(p-1),\qquad \B_0=\F_p,\qquad \B=\F_{p^2},
\]
let \(H\le\B^\times\) have order \(n\), choose a generator \(\theta\) of
\(\B^\times\), and set \(D=\theta H\).  Choose even \(a=a_p\) with
\(a/n\to\alpha\), let \(w\) be the largest even integer not exceeding
\(\log_{|\B|}\binom na\), and put \(k=a-w-1\).  Then
\[
 \log|\B|=O(\log n)=o(n),\quad p>a,\quad
 w=\Theta(n/\log n),\quad 1\le\barN_1<|\B|^2=e^{o(n)}.
 \tag{5.2}\label{eq:counterexample-scale}
\]
For some \(z\in\B^w\),
\[
 |\mathcal F_z|\ge \binom{n/2}{a/2}p^{-w/2}
 =\exp\left(\left(\frac14h(\alpha)+o(1)\right)n\right),
 \qquad h(x)=-x\log x-(1-x)\log(1-x).
 \tag{5.3}\label{eq:counterexample-fiber}
\]
The quantity before pigeonholing is the natural average
\[
 \barN_{\rm sq}=\binom{n/2}{a/2}|\B_0|^{-w/2}
\]
of the complete-square quotient profile.  Thus every exhaustive profile
atlas containing this quotient has
\(\Eprof_n(a)\ge\barN_{\rm sq}
=\exp((h(\alpha)/4+o(1))n)\), whereas
\(\barN_1=e^{o(n)}\); in particular the row is not identity-dominant
there.
There is \(\sigma=\sigma(\alpha)>0\) such that
\(\Delta(\mathcal F_z)\le e^{-\sigma n}\), and every logarithmic
\(r=r_n\to\infty\) satisfies
\[
 \liminf_{n\to\infty}\frac1{nr}
       \log\widetilde\Gsid_{r,\sigma}
 \ge\frac14h(\alpha)>0.
 \tag{5.4}\label{eq:counterexample-sidon}
\]
Finally, every finite extension \(\F/\B\) with
\[
 |\F|-n>k\binom{|\mathcal F_z|}{2}
 \tag{5.5}\label{eq:counterexample-extension}
\]
satisfies
\[
 B_{\RS_\F(D,k)}^{\rm MCA}(a)\ge|\mathcal F_z|
 \ge\exp\left(\left(\frac14h(\alpha)+o(1)\right)n\right).
 \tag{5.6}\label{eq:counterexample-mca}
\]
Such extensions exist with degree \(O(n/\log n)\).  Hence smoothness and
\(\log|\B|=o(n)\) do not imply an identity-normalized Sidon payment or an
identity-scale MCA numerator bound.  Indeed,
\[
        B_{\RS_\F(D,k)}^{\rm MCA}(a)=\exp(\Omega(n))
        \qquad\text{while}\qquad
        \barN_1=\exp(o(n)).
\]
\end{theorem}

\begin{proof}
The cyclic group \(\B^\times\) contains \(H\), and \(H^2\) has order
\(p-1\); hence \(H^2=\B_0^\times\).  Thus
\(D^2=\theta^2\B_0^\times\), and squaring is two-to-one on \(D\).  The
roots of \(1+X^2\) have order four and lie in \(H\), whereas the proper
coset \(D=\theta H\) is disjoint from \(H\).  Since
\(a/n\to\alpha<1/2\), eventually \(a<p\).  Stirling's formula and the
even rounding give
\(w=nh(\alpha)/(2\log p)+o(n/\log n)\) and
\(0\le\log_{|\B|}\binom na-w<2\), proving
\eqref{eq:counterexample-scale}.

For \(E\in\binom{D^2}{a/2}\), let
\(S_E=\{x\in D:x^2\in E\}\).  Its locator is
\[
 Q_{S_E}(X)=\prod_{y\in E}(X^2-y)
 =X^a-e_1(E)X^{a-2}+e_2(E)X^{a-4}-\cdots .
\]
The odd-gap coefficients vanish, while the coefficient in gap \(2j\)
lies in \(\theta^{2j}\B_0\).  Hence these supports occupy at most
\(p^{w/2}\) prefixes.  Pigeonholing gives the first inequality in
\eqref{eq:counterexample-fiber}; moreover
\(\log\binom{n/2}{a/2}=\tfrac12nh(\alpha)+o(n)\) and
\((w/2)\log p=\tfrac14nh(\alpha)+o(n)\), proving the exponent.

Let \(A\) be the constant in
\cref{cor:quantitative-boolean-energy}, and choose
\(0<\sigma<h(\alpha)/(8A)\).  If
\(\Delta(\mathcal F_z)>e^{-\sigma n}\), that corollary with
\(K=e^{\sigma n}\) gives \(|\mathcal F_z|\le e^{A\sigma n}\),
contradicting \eqref{eq:counterexample-fiber}.  Keeping this single
summand in \eqref{eq:unrestricted-sidon-moment} gives
\[
 \log\widetilde\Gsid_{r,\sigma}
 \ge-w\log|\B|+r(\log|\mathcal F_z|-\log\barN_1)
 =-nh(\alpha)+o(n)+r\left(\frac14h(\alpha)n+o(n)\right),
\]
which proves \eqref{eq:counterexample-sidon}.  Because \(w<a<p\),
Newton identities give an invertible triangular change between these
locator coefficients and the first \(w\) power sums.

Set \(U_z=X^a+z_1X^{a-1}+\cdots+z_wX^{a-w}\) and
\(P_S=U_z-Q_S\).  Then \(\deg P_S\le k\), and for distinct
\(S,T\in\mathcal F_z\), the nonzero polynomial \(Q_S-Q_T\) has degree at
most \(k\).  Condition \eqref{eq:counterexample-extension} leaves
\(\lambda\in\F\setminus D\) at which all \(Q_S(\lambda)\) are distinct.
On \(D\) put \(r_1(x)=U_z(x)/(x-\lambda)\) and
\(r_2(x)=-1/(x-\lambda)\).  For
\(\gamma_S=P_S(\lambda)\) and \(x\in S\),
\[
 r_1(x)+\gamma_Sr_2(x)
 =\frac{P_S(x)-P_S(\lambda)}{x-\lambda},
\]
which is a polynomial of degree less than \(k\).  The slopes are distinct.
Moreover \(r_2\) cannot agree with a degree-less-than-\(k\) polynomial
on \(k+1\) points, since then \((X-\lambda)q(X)+1\) would have
\(k+1\) roots, degree at most \(k\), and value one at \(\lambda\).
Thus all \(\gamma_S\) are MCA-bad.  Finally
\(\log|\mathcal F_z|=O(n)\) and \(\log|\B|=\Theta(\log n)\), so an
extension of degree \(O(n/\log n)\) suffices.
\end{proof}

\begin{corollary}[Circle obstruction]\label{cor:circle-obstruction}
In \cref{thm:polynomial-obstruction}, put
\(u=(k-1)/2=(a-w-2)/2\) and take the inverse stereographic image of \(D\)
on \(x^2+y^2=1\).  Let \(C_{\rm circ}(u)\) be the evaluation, over
\(\F\), of
\[
 \{f_0(x)+yf_1(x):\deg f_0\le u,\ \deg f_1\le u-1\}
\]
on that circle set.  Then \(C_{\rm circ}(u)\) is diagonally equivalent to
\(\RS_\F(D,k)\); hence it has the same MCA-bad slopes and the same lower
bound.
\end{corollary}

\begin{proof}
The inverse parametrization is
\(x=(1-t^2)/(1+t^2)\), \(y=2t/(1+t^2)\).  Its denominator is nonzero on
\(D\) by the order-four argument above.  Modulo \(x^2+y^2-1\), the
degree-\(u\) space is
\(\{f_0(x)+yf_1(x):\deg f_0\le u,\ \deg f_1\le u-1\}\).  After
substitution and multiplication by \((1+t^2)^u\), it becomes exactly the
polynomials in \(t\) of degree at most \(2u\): substitution is injective
in the localized coordinate ring, and both spaces have dimension
\(2u+1\).  This coordinatewise nonzero multiplier preserves all support
equalities and slopes.  Finally, \(a,w\) are even, so
\(k=a-w-1=2u+1\).
\end{proof}

\begin{remark}[Stabilizer deletion is insufficient]
Fix \(x_0\in D\) and use
\(S_E=\{x_0\}\cup\{x\in D:x^2\in E\}\), where
\(E\in\binom{D^2\setminus\{x_0^2\}}{a/2}\).  At support size \(a+1\)
and depth \(w\), some prefix \(z\) has a subfamily
\[
 \mathcal F_z^{\rm ap}\subseteq\Phi_w^{-1}(z),\qquad
 |\mathcal F_z^{\rm ap}|
 \ge\binom{n/2-1}{a/2}p^{-w/2}
 =\exp\left(\left(\frac14h(\alpha)+o(1)\right)n\right).
\]
Here every member has trivial multiplicative stabilizer.  Indeed a
stabilizer of a nonempty subset of \(D\) lies in
\(DD^{-1}=H\), and each \(S_E\) has the unique incomplete orbit
\(\{x_0,-x_0\}\) under \(-1\).  Its stabilizer must preserve this pair and
is therefore \(\pm1\), while \(-1\) does not preserve \(S_E\).  Moreover
\(\binom n{a+1}|\B|^{-w}=e^{o(n)}\), so
\cref{cor:quantitative-boolean-energy} gives
\(\Delta(\mathcal F_z^{\rm ap})\le e^{-\sigma n}\) for some fixed
\(\sigma>0\).  This is a Sidon-heavy fiber inside the
trivial-\(H\)-stabilizer residual.  It is a planted remainder profile, so
aperiodicity alone does not make the unrestricted identity ledger valid.
\end{remark}

\section{MCA compiler}

\begin{lemma}[Primitive add-back]\label{lem:addback}
For a closed-ledger row, let \(N_\lambda(s)\) be the residual prefix
fibers in the primitive profile \(\lambda\), after all algebraic cells have
been removed.  Then
\[
 \sum_{\lambda\ {\rm primitive}}\max_sN_\lambda(s)
 \le\exp(o(n))\Eprof_n(a_n).
\]
\end{lemma}

\begin{proof}
Primitive leaves satisfy \cref{thm:primitive-q}, so their maximal fibers are
\(\exp(o(n))M_\lambda/L_\lambda\).  The residual-to-full comparison in
\cref{def:closed-ledger} bounds this by
\(\exp(o(n))\barN_\lambda\), which occurs in
\eqref{eq:profile-envelope}.  Summing over the subexponential first-match
profile family preserves the same scale.
\end{proof}

\begin{lemma}[Q pays shift pairs]\label{lem:q-sp}
If \(\sum_sN(s)=M\), \(\barN=M/L\), and \(\max_sN(s)\le \kappa\barN\), then \(M^{-1}\sum_sN(s)^2\le\kappa\barN\).
\end{lemma}

\begin{proof}
\(\sum_sN(s)^2\le(\max_sN(s))\sum_sN(s)\le\kappa\barN M\).
\end{proof}

\begin{remark}[Q and SP do not count rays]\label{rem:q-sp-no-ray}
The preceding lemma is a support-pair identity, not a ray compiler.  Keeping
the same fibers \(N(s)\), an abstract incidence can assign distinct slopes
to selected supports or assign one slope to all of them without changing
\(\max_sN(s)\) or \(\sum_sN(s)^2\).  An RS numerator bound therefore
requires the additional algebraic projection estimate RC.
\end{remark}

\begin{theorem}[Closed-ledger upper bound]\label{thm:upper}
If a smooth RS row sequence is closed-ledger and satisfies RC at \(a_n\),
then
\[
        B_{C_n}^{\rm MCA}(a_n)
        \le \exp(o(n))\Eprof_n(a_n).
\]
\end{theorem}

\begin{proof}
The nonprimitive cells are paid by definition.  The primitive prefix leaves
obey Q by \cref{thm:primitive-q} and add back by \cref{lem:addback}; their
support-pair second moment is bounded by \cref{lem:q-sp}.  By
\cref{rem:q-sp-no-ray}, this does not itself bound slopes.  RC supplies the
required slope-image estimate, and first-match disjointization sums the
disjoint slope classes.
\end{proof}

\section{Lower side and frontier}

\begin{proposition}[Collision-aware identity lower bound]
\label{prop:collision-aware-lower}
Let \(D\subseteq\B\subseteq\F\), let \(C=\RS_\F(D,k)\),
\(q=|\F|>n\), and \(k+1\le m\le n\).  Put
\[
 w=m-k-1,
 \qquad L=\left\lceil\binom nm|\B|^{-w}\right\rceil.
\]
Then
\[
 B_C^{\rm MCA}(m)\ge
 \left\lceil\frac{L(q-n)}{q-n+k(L-1)}\right\rceil.
 \tag{7.1}\label{eq:collision-aware-lower}
\]
\end{proposition}

\begin{proof}
Pigeonhole the first \(w\) locator coefficients to obtain \(L\) supports
with common prefix \(z\).  If
\(Q_S=\prod_{x\in S}(X-x)\) and
\(U_z=X^m+z_1X^{m-1}+\cdots+z_wX^{m-w}\), then
\(P_S=U_z-Q_S\) has degree at most \(k\).  At a pole
\(\lambda\in\F\setminus D\), the simple-pole line used in the proof of
\cref{thm:polynomial-obstruction} makes every distinct value
\(P_S(\lambda)\) an MCA-bad slope.

For \(S\ne T\), equality \(P_S(\lambda)=P_T(\lambda)\) occurs at at most
\(k\) poles.  Averaging over the \(q-n\) poles gives one pole with at most
\(k\binom L2/(q-n)\) colliding pairs.  If its value multiplicities are
\(m_1,\ldots,m_M\), then
\[
 \sum_i m_i^2\le L+\frac{kL(L-1)}{q-n}.
\]
Cauchy--Schwarz, \(L^2\le M\sum_i m_i^2\), proves
\eqref{eq:collision-aware-lower}.  This is the quantitative form of the
prefix-pole construction in \cite{Cho26CapV13,Cho26Grande}.
\end{proof}

\begin{proof}[Proof of \cref{thm:frontier}]
Condition \textup{(i)} and \cref{thm:upper} make \(a_{+,n}\) safe,
while condition \textup{(ii)} makes \(a_{-,n}\) unsafe.  Since
\(B_{C_n}^{\rm MCA}(a)\) is nonincreasing in \(a\), the first safe
agreement satisfies \(a_{-,n}<a_n^*\le a_{+,n}\), proving
\eqref{eq:threshold-bracket}.  If the endpoints are \(o(n)\) apart around
\(k_n+1+g_nn\), division by \(n\) gives
\(\delta_n^*=1-\rho_n-g_n+o(1)\).  The displayed identity choice of
\(U_n(a)\) is \cref{prop:collision-aware-lower}.

For \(a_n=k_n+1+gn+O(1)\), Stirling's formula gives
\[
        \log_2\barN_{n,a_n}
        =nF_n(g)+o(n).
\]
Thus \(g_{T,n}\) is the exponential-scale right boundary at which the
identity average has target rate \(\tau_n\).  The separate endpoint
hypothesis converts this exponential comparison into the exact safe and
unsafe inequalities, accounting for the tangent term, compiler loss,
integer rounding, and pole collisions.  The preceding bracket gives the
target-adjusted formula.  Finally, if \(\tau_n=o(1)\), concavity of
\(F_n\), together with \(F_n(0)=H_2(\rho_n)\) bounded away from zero,
shows that its right \(\tau_n\)-level boundary differs by \(o(1)\) from
its right zero-level boundary \(g^*(\rho_n,\beta_n)\).  This proves
\eqref{eq:identity-frontier}.
\end{proof}

\begin{remark}[Scope of the main statement]
\Cref{thm:polynomial-obstruction} shows that smooth/circle geometry and
\(\log|\B_n|=o(n)\) do not imply the identity-scale upper bound: a
degree-two quotient over a proper subfield already has an exponentially
larger natural prefix budget.  Thus the unconditional conclusions here are
the obstruction and the high-energy elimination.  The frontier formula is
the identity-dominant specialization of the profile-envelope compiler.  For
unrestricted rows, the remaining requirements are exhaustive
quotient/remainder routing with residual-to-full profile comparison, a
Sidon payment on the resulting primitive residual, and a
higher-dimensional transverse-ray compiler.
Because the constructed depth satisfies \(w/n\to0\), the example alone
does not disprove the coarse \(o(1)\)-radius formula for subexponential
targets; it disproves the identity-scale numerator bound and hence any
unconditional derivation of that formula from the identity prefix alone.
\end{remark}

\enlargethispage{4\baselineskip}
\begin{thebibliography}{GMRSZ20}

\bibitem[BS94]{BalogSzemeredi}
A. Balog and E. Szemeredi,
\newblock A statistical theorem of set addition,
\newblock \emph{Combinatorica} 14 (1994), 263--268.

\bibitem[Cho26a]{Cho26CapV13}
P. Chojecki,
\newblock \emph{Universal Field-Size Caps and a Two-Sided Sandwich for
Mutual Correlated Agreement on Smooth Reed--Solomon Domains},
\newblock manuscript, July 2026.

\bibitem[Cho26b]{Cho26Grande}
P. Chojecki,
\newblock \emph{Proof-Audited and Expanded Bounds and Reductions for
RS--MCA},
\newblock manuscript, July 2026.

\bibitem[GMRSZ20]{GMRSZ2020}
B. Green, D. Matolcsi, I. Z. Ruzsa, G. Shakan, and D. Zhelezov,
\newblock A weighted Prekopa--Leindler inequality and sumsets with quasicubes,
\newblock arXiv:2003.04077, 2020.

\bibitem[Gow98]{Gowers1998}
W. T. Gowers,
\newblock A new proof of Szemeredi's theorem for arithmetic progressions of length four,
\newblock \emph{Geom. Funct. Anal.} 8 (1998), 529--551.

\bibitem[MRSZ20]{MRSZ2020}
D. Matolcsi, I. Z. Ruzsa, G. Shakan, and D. Zhelezov,
\newblock An analytic approach to cardinalities of sumsets,
\newblock arXiv:2003.04075, 2020.

\end{thebibliography}

\end{document}
